\begin{lemma}
\label{lemma-tor-amplitude}
Let $(X, \mathcal{O}_X)$ be a ringed space.
Let $E$ be an object of $D(\mathcal{O}_X)$.
Let $a, b \in \mathbf{Z}$ with $a \leq b$. The following are equivalent
\begin{enumerate}
\item $E$ has tor-amplitude in $[a, b]$.
\item $E$ is represented by a complex
$\mathcal{E}^\bullet$ of flat $\mathcal{O}_X$-modules with
$\mathcal{E}^i = 0$ for $i \not \in [a, b]$.
\end{enumerate}
\end{lemma}

\begin{proof}
If (2) holds, then we may compute
$E \otimes_{\mathcal{O}_X}^\mathbf{L} \mathcal{F} =
\mathcal{E}^\bullet \otimes_{\mathcal{O}_X} \mathcal{F}$
and it is clear that (1) holds.

\medskip\noindent
Assume that (1) holds. We may represent $E$ by a bounded above complex
of flat $\mathcal{O}_X$-modules $\mathcal{K}^\bullet$, see
Section \ref{section-flat}.
Let $n$ be the largest integer such that $\mathcal{K}^n \not = 0$.
If $n > b$, then $\mathcal{K}^{n - 1} \to \mathcal{K}^n$ is surjective as
$H^n(\mathcal{K}^\bullet) = 0$. As $\mathcal{K}^n$ is flat we see that
$\Ker(\mathcal{K}^{n - 1} \to \mathcal{K}^n)$ is flat
(Modules, Lemma \ref{modules-lemma-flat-ses}).
Hence we may replace $\mathcal{K}^\bullet$ by
$\tau_{\leq n - 1}\mathcal{K}^\bullet$. Thus, by induction on $n$, we
reduce to the case that $K^\bullet$ is a complex of flat
$\mathcal{O}_X$-modules with $\mathcal{K}^i = 0$ for $i > b$.

\medskip\noindent
Set $\mathcal{E}^\bullet = \tau_{\geq a}\mathcal{K}^\bullet$.
Everything is clear except that $\mathcal{E}^a$ is flat
which follows immediately from Lemma \ref{lemma-last-one-flat}
and the definitions.
\end{proof}
